\documentclass[11pt]{article}

\usepackage[T1]{fontenc}
\usepackage[utf8]{inputenc}
\usepackage{lmodern}
\usepackage{microtype}
\usepackage[margin=1in]{geometry}
\usepackage{amsmath,amssymb,amsthm,mathtools}
\usepackage{booktabs,longtable,array}
\usepackage{enumitem}
\usepackage{xcolor}
\usepackage[hidelinks]{hyperref}
\hypersetup{
  pdftitle={From Unequal Beginnings to Social Starting Capital},
  pdfauthor={Housam Kak},
  pdfsubject={A normative and budget-closed framework for common inheritance, child development, and protected capital},
  pdfkeywords={social starting capital, common inheritance, child accounts, inequality, intergenerational justice}
}
\Urlmuskip=0mu plus 1mu\relax

\newtheorem{definition}{Definition}[section]
\newtheorem{proposition}[definition]{Proposition}
\newtheorem{corollary}[definition]{Corollary}
\newtheorem{remark}[definition]{Remark}
\newtheorem{assumption}[definition]{Assumption}

\newcommand{\R}{\mathbb{R}}
\newcommand{\one}{\mathbf{1}}
\newcommand{\pos}[1]{\left[#1\right]_{+}}
\newcommand{\clamp}[1]{\operatorname{clamp}_{[0,1]}\!\left(#1\right)}

\title{From Unequal Beginnings to Social Starting Capital\\
\large A Normative and Budget-Closed Framework for Common Inheritance,\\
Child Development, and Protected Capital}
\author{Housam Kak}
\date{October 2026}

\begin{document}

\maketitle

\begin{abstract}
People do not choose the families, assets, legal systems, public institutions, or
historical circumstances into which they are born. Yet those conditions shape their
capabilities long before individual effort can plausibly explain the resulting
differences. This paper argues that part of a society's productive inheritance---the
value attributable to nature, accumulated knowledge, public institutions, social
cooperation, and historically created scarcity rents---may legitimately finance a
material starting stake for every child. It then develops the Social Starting Capital
Framework (SSCF), a policy architecture with three distinct functions: a universal
child-development floor, supplementary household support, and an independently owned
birthright capital account.

The formal model is designed to be budget-closed. Every newborn receives an equal,
strictly positive statutory base grant. Family size affects supplementary household
support, not the child's independent capital right. All discretionary bonuses and
progressive top-ups are normalized to explicit pools, so realized expenditure cannot
exceed the authorized budget when births or household composition differ from
forecasts. The model distinguishes maximum spending capacity from actual spending,
specifies the split between current revenue and fund withdrawals, includes withdrawals
in child-account dynamics, and states an exact condition for preserving the social
wealth fund in real per-capita terms. A fiscal backstop makes the claimed minimum a
genuine guarantee rather than a residual aspiration. A calibration and demographic
stress test illustrate the system.

The proposal is neither a theory of equal outcomes nor a substitute for public
services, private saving, or claims of historical rectification. It is a pre-
distributional institution: each generation inherits a protected material foothold in
a social world it did not create but must be able to enter as an agent.
\end{abstract}

\newpage
\tableofcontents
\newpage

\part{The Normative Case}

\section{Introduction}

The conditions of birth are unchosen, but they are not inconsequential. A child may be
born into secure housing, educated networks, inherited assets, political stability,
and ready access to credit. Another may be born into debt, displacement, dangerous
housing, weak public services, or legal insecurity. Neither child authored that initial
position. By the time conventional accounts of merit begin measuring effort, the two
children may already differ in health, vocabulary, confidence, time horizon, and the
range of choices they can safely make.

This fact does not imply that every later inequality is unjust. Differences can arise
from effort, preference, innovation, luck, care obligations, and voluntary exchange.
The narrower claim defended here is that a political order cannot treat radically
unequal starting power as morally irrelevant and then describe the resulting
competition as equal opportunity. Formal permission is not the same as effective
freedom. A right to study, move, refuse an abusive employer, or start an enterprise is
thin when a person lacks the health, time, security, and assets required to exercise it.

The empirical background is severe. The \emph{World Inequality Report 2022} estimates
that the richest tenth of the global population receives roughly 52 percent of global
income and owns roughly 76 percent of global wealth, while the poorest half receives
8.5 percent of income and owns about 2 percent of wealth \cite{wir2022}. Land---a
finite and politically constituted asset---is also highly concentrated; in the most
unequal regions, the top tenth of landowners controls as much as three quarters of
agricultural land \cite{bauluz2020}. Measures of intergenerational mobility differ by
country and by whether one studies income, education, absolute mobility, or relative
mobility, but family background remains a powerful predictor of later outcomes
\cite{vanderweide2024}.

The policy response developed in this paper begins from a concept of \emph{common
inheritance}. Each generation receives natural resources, scientific knowledge,
language, infrastructure, legal forms, monetary institutions, and productive
organizations that it did not build from nothing. Those inheritances do not make all
wealth collective. They do, however, weaken the claim that the entire social product is
morally exhausted by present private title.

The Social Starting Capital Framework (SSCF) turns that moral observation into an
institution with three separate layers:

\begin{enumerate}[label=\arabic*.]
  \item a universal child-development floor, delivered primarily through services or
  guaranteed provision;
  \item supplementary household support, recognizing that care occurs in families and
  has both fixed and variable costs; and
  \item a protected capital account owned by the child, containing an equal statutory
  base grant and, where authorized, budget-normalized additions.
\end{enumerate}

Separating these layers is essential. Services cannot provide the bargaining power of
an owned asset. A capital account cannot replace nutrition, healthcare, safe housing,
or schools. Household assistance is necessary, but a payment controlled entirely by
parents is not the same institution as capital legally owned by the child.

The remainder of Part I develops the moral argument and its limits. Part II provides a
complete accounting model. Part III gives a calibration, stress test, governance rules,
and limitations.

\section{Unequal beginnings and effective freedom}

The metaphor of the birth lottery is useful because it identifies a boundary of
responsibility. No one chooses their parents, citizenship, childhood neighborhood,
inherited health, first language, or initial balance sheet. Rawls treats natural and
social starting positions as morally arbitrary for the purpose of choosing principles
for the basic structure \cite{rawls1971}. Sen shifts attention from formal holdings to
capabilities: what a person is substantively able to be and do \cite{sen1999}. Van
Parijs similarly distinguishes formal freedom from the resources required for real
freedom \cite{vanparijs1995}.

Capital matters within this framework for reasons different from income. Income meets
recurrent needs. Public services create capabilities and insure risks. Capital adds a
stock of option value. It can finance relocation, education, tools, a housing deposit,
or a period of search. It can make refusal possible. An individual with a modest asset
buffer can reject a predatory loan or dangerous job in a way that an otherwise equal
individual living at the edge of subsistence cannot.

This paper therefore uses \emph{starting capital} in the financial and institutional
sense. It should not be confused with the sociological concept of \emph{social capital},
which normally refers to relationships, networks, trust, and norms. Networks are part
of unequal beginnings, but the SSCF is a proposal for social starting \emph{capital}:
a publicly constituted material stake.

\section{Property, history, and common inheritance}

Modern holdings did not emerge on an empty field. They are the result of labor and
exchange, but also of inheritance, enclosure, conquest, slavery, colonial extraction,
state-building, technological accumulation, and legal decisions about which claims the
state will enforce. Recognizing this history does not establish that every present title
is invalid. It establishes that present title cannot be the only premise in a theory of
distribution.

Locke's labor theory of appropriation remains important because effort and personal
projects deserve protection. Yet Lockean appropriation is constrained by the proviso
that enough and as good remain for others \cite{locke1689}. Once land and resources are
scarce and legally enclosed, latecomers encounter neither open nature nor a blank slate.
Rousseau emphasizes that civil inequality is institutional, not merely a reflection of
natural difference \cite{rousseau1755}. Paine argues that private cultivation can be
recognized while those excluded from the earth's original common condition receive
compensation \cite{paine1797}. George makes a related distinction between productive
effort and land rent generated by scarcity and surrounding social development
\cite{george1879}.

Marx's account of primitive accumulation and Polanyi's account of land, labor, and
money as politically constructed commodities place market exchange within a prior
history of institutional change \cite{marx1867,polanyi1944}. Their analyses need not
be accepted in full to recognize that markets depend on law, public authority,
infrastructure, money, and accumulated knowledge.

Nozick's entitlement theory offers the strongest warning against treating persons and
their holdings as instruments of a preferred pattern. Just holdings acquired and
transferred legitimately deserve protection \cite{nozick1974}. But the historical
principle also makes acquisition and rectification unavoidable questions. A society
with layered dispossession, inherited legal privilege, and uncompensated externalities
cannot simply assume the clean history that gives the entitlement argument its greatest
force.

The conclusion is limited but significant: a defined portion of economic flows can be
treated as socially claimable without treating all property or production as collective.
The relevant base is not ``everything the state can tax.'' It is a transparent set of
revenues connected to common assets, inherited transfers, publicly created value,
external costs, and collectively owned capital.

This conclusion is strengthened by two further observations. First, wealth can compound
across generations and become increasingly important to stratification when returns on
owned assets outpace the growth of ordinary incomes \cite{piketty2014}. Second,
recognizing a common claim does not require a choice between atomized private ownership
and unrestricted central control. Ostrom's work demonstrates that shared resources can
be governed through durable, rule-bound institutions under the right conditions
\cite{ostrom1990}. The design question is therefore not simply public versus private,
but which rights, duties, boundaries, and monitoring arrangements fit the resource.

\section{Three arguments that must remain distinct}

Several arguments for redistribution are often blended together. They should be
separated because they imply different allocation rules.

\subsection{Common inheritance}

If every member of a generation is a co-beneficiary of natural, institutional, and
knowledge inheritances, the common-inheritance argument supports a universal claim.
The equal base grant in this paper rests primarily on that argument.

\subsection{Capability and priority}

Equal nominal grants do not equalize substantive freedom. Disability, regional prices,
discrimination, care burdens, and household poverty affect what the same amount can do.
Capability and prioritarian arguments therefore support progressive or needs-based
top-ups. Those additions should be explicit and separately budgeted rather than hidden
inside the universal claim.

\subsection{Rectification}

Historical injustice can generate claims for groups or persons who suffered or inherit
the consequences of a specific wrong. Such claims may be larger, differently funded,
or differently governed than a universal starting stake. A universal SSCF must not be
presented as extinguishing claims for reparations, restitution, indigenous sovereignty,
or other forms of rectification.

This separation improves both philosophy and policy. Universality, progressivity, and
rectification can coexist, but one does not silently substitute for another.

\section{Intellectual and institutional lineage}

The proposal belongs to a long family of arguments. Paine proposed a national fund that
would pay a capital sum at adulthood \cite{paine1797}. More recent proposals include
stakeholder grants, minimum inheritance, and asset-building policies
\cite{ackerman1999,atkinson2015,sherraden1991}. Hamilton and Darity's baby-bonds work
shows how progressively funded child accounts can address entrenched wealth gaps
\cite{hamilton2010}.

Public wealth funds provide a second lineage. Norway's Government Pension Fund turns
petroleum revenue into long-term public saving and explicitly links its management to
intergenerational considerations \cite{norwayfund}. These examples do not by themselves
justify the SSCF, nor are they institutionally identical to it. They establish that
public ownership, professional asset management, and intergenerational distribution are
practical institutional categories rather than conceptual impossibilities.

The SSCF's contribution is architectural. It integrates present capabilities, family
care, individually owned capital, explicit fiscal priority, and a fund rule in one
budget-closed model. Its novelty claim is therefore about the combination and the
formal separation of functions, not the invention of child accounts or public funds.

\section{Normative principles}

The framework is governed by seven principles.

\begin{description}[style=nextline]
  \item[P1. Common inheritance.] A defined part of social wealth is attributable to
  nature, inherited institutions, public action, and cumulative cooperation rather than
  current private effort alone.

  \item[P2. Independent child standing.] The child is a rights-bearing unit. Every
  eligible newborn receives an equal, strictly positive base capital grant that is not
  reduced by sibling count, parental income, or module completion.

  \item[P3. Capability before discretion.] Essential childhood development has fiscal
  priority over optional distributions.

  \item[P4. Family agency without substitution.] Families receive room and support to
  care, but household transfers neither replace the child-development floor nor absorb
  the child's capital right.

  \item[P5. Exact budget closure.] Every variable expenditure is assigned to a defined
  pool. Forecast errors may change bonuses, but cannot create unauthorized expenditure.

  \item[P6. Intergenerational reciprocity.] Fund policy is evaluated in real per-capita
  terms, not merely by nominal non-decline.

  \item[P7. Non-forfeiture and proportionate protection.] Education may affect the form
  and timing of early access, but never ownership. Restrictions expire and are subject
  to due process.
\end{description}

\part{The Formal Framework}

\section{Timing, populations, and notation}

Time is discrete and indexed by fiscal years $t=0,1,2,\ldots$. All transfers for birth
cohort $t$ are finalized after the eligible births in that cohort have been registered.
A provisional account may be opened at birth, but the final annual bonus is credited
only after the realized cohort is known. This convention eliminates the use of future
or unobserved population values in current grants.

Let:

\begin{longtable}{>{\raggedright\arraybackslash}p{0.18\textwidth}p{0.72\textwidth}}
\toprule
Symbol & Meaning \\
\midrule
$N_t$ & reference population at the start of year $t$ \\
$\mathcal F_t$ & set of eligible households during year $t$ \\
$K_{f,t}$ & number of eligible dependent children in household $f$ \\
$\mathcal C_t$ & set of all eligible dependent children \\
$\mathcal B_t$ & finalized set of eligible newborns in cohort $t$ \\
$B_t$ & realized births, $B_t=|\mathcal B_t|$ \\
$Y_t$ & net, non-overlapping socially claimable inflow \\
$W_t$ & market value of the social wealth fund at the start of year $t$ \\
$r_t$ & net nominal return on the fund during year $t$ \\
$P_t$ & general price index, used only for real-value comparisons \\
$U_t$ & expenditure on the universal child-development floor \\
$G_t^{\min}$ & aggregate statutory newborn base grants \\
$H_t$ & actual supplementary household allocation \\
$G_t^{+}$ & cohort bonus above the statutory base \\
$T_t$ & progressive top-up pool for existing child accounts \\
$V_t$ & optional adult or transition distribution \\
$D_t$ & actual total program distribution \\
$D_t^Y,D_t^W$ & actual distribution funded from current resources and the fund \\
$Z_t$ & transparent general-revenue backstop for mandatory child claims \\
$A_{c,s}$ & total account assets owned by child $c$ at date $s$ \\
\bottomrule
\end{longtable}

We assume $N_t>0$, $P_t>0$, $W_t\ge 0$, and $B_t\ge 0$. Every formula that divides by
$B_t$ includes a separate zero-birth convention.

\section{The revenue base and social wealth fund}

\begin{definition}[Socially claimable inflow]
The net inflow dedicated to the framework is
\begin{equation}
  Y_t=Y_t^{\mathrm{res}}+Y_t^{\mathrm{land}}+Y_t^{\mathrm{inh}}
      +Y_t^{\mathrm{pub}}+Y_t^{\mathrm{eco}}+Y_t^{\mathrm{other}},
\end{equation}
where the terms denote, respectively, net resource rents, land-value revenue, dedicated
inheritance and gift taxation, distributable profit from public assets, ecological
charges, and other legislatively designated revenue.
\end{definition}

The components must be measured net of collection costs, rebates, and prior legal
claims. They must also be non-overlapping. In particular, $Y_t^{\mathrm{pub}}$ excludes
returns already counted in $r_tW_t$, and ecological charges exclude resource royalties
already counted as resource rents. This restriction prevents the same economic flow
from financing the model twice.

Let $\phi_t\in[0,1]$ be the minimum saving share of $Y_t$. Let $\widetilde W_t$ be a
prudential reference value known at the start of the year, such as the lesser of current
fund value and a trailing multi-year average. With draw-rate parameter
$\gamma_t\in[0,1)$, the maximum authorized fund draw is
\begin{equation}
  \overline d_t=\min\{\gamma_t\widetilde W_t,W_t\}.
\end{equation}
The use of a predetermined reference value avoids making current entitlements depend on
an end-of-year return that is not yet known.

Before the statutory backstop is applied, maximum current-year spending capacity is
\begin{equation}
  C_t^0=(1-\phi_t)Y_t+\overline d_t.
\end{equation}
This is capacity, not actual expenditure and not an instruction to take the full draw.

\section{Mandatory child claims and the fiscal backstop}

Let $u_{c,t}\ge0$ be the cost of the guaranteed development package for child $c$. It
may reflect disability, regional cost, or age while preserving a common service
standard. Then
\begin{equation}
  U_t=\sum_{c\in\mathcal C_t}u_{c,t}.
\end{equation}
The package can include nutrition, preventive healthcare, schooling, safe transport,
digital access, and housing-related support. It should be measured as the incremental
cost attributable to the SSCF or explicitly identify existing expenditure reclassified
under it; otherwise the fiscal cost would be overstated or double-counted.

Let $\underline g_t>0$ be the indexed statutory base grant for each eligible newborn.
The aggregate base obligation is
\begin{equation}
  G_t^{\min}=B_t\underline g_t.
\end{equation}
The mandatory claim is
\begin{equation}
  M_t=U_t+G_t^{\min}.
\end{equation}

A right cannot honestly be called guaranteed if it disappears whenever residual funds
are exhausted. The framework therefore makes any shortfall visible:
\begin{equation}
  Z_t=\pos{M_t-C_t^0}.
\end{equation}
Here $Z_t$ is a transfer from ordinary public revenue authorized by the statute creating
the SSCF. It may be financed through a standing appropriation or an automatic surcharge.
It may not be financed by breaching the fund draw cap. Total authorized capacity is
therefore
\begin{equation}
  C_t=C_t^0+Z_t,
\end{equation}
and $M_t\le C_t$ by construction.

\begin{remark}
The backstop is not free money. It states who bears the risk when dedicated revenues
fall short. Without it, either the minimum is not a right or the model contains an
unreported deficit. A legislature may set a modest $\underline g_t$ to make activation
rare, but it cannot claim an unconditional positive grant while assigning all revenue
risk to the child.
\end{remark}

\section{Supplementary household allocation}

Households face fixed costs and economies of scale. Let the desired supplemental amount
for household $f$ be
\begin{equation}
  \widehat H_{f,t}=h_t+s_t\omega(K_{f,t}),
  \qquad h_t,s_t\ge0,
\end{equation}
where
\begin{equation}
  \omega(K)=\sum_{j=1}^{K}w_j,
  \qquad
  w_j=
  \begin{cases}
    1, & j\le K^*,\\
    \lambda^{j-K^*}, & j>K^*,
  \end{cases}
\end{equation}
with $K^*\in\{1,2,\ldots\}$ and $0<\lambda<1$. The taper applies only to this
supplementary household layer. It never changes $\underline g_t$.

Let
\begin{equation}
  \widehat H_t=\sum_{f\in\mathcal F_t}\widehat H_{f,t}
\end{equation}
be desired aggregate household support. A share $\beta_H\in[0,1]$ of capacity remaining
after mandatory claims is reserved for this layer:
\begin{equation}
  \overline H_t=\beta_H(C_t-M_t).
\end{equation}
Define
\begin{equation}
  \xi_t=
  \begin{cases}
    \min\{1,\overline H_t/\widehat H_t\}, & \widehat H_t>0,\\
    0, & \widehat H_t=0,
  \end{cases}
\end{equation}
and set
\begin{equation}
  H_{f,t}=\xi_t\widehat H_{f,t},
  \qquad H_t=\sum_{f\in\mathcal F_t}H_{f,t}.
\end{equation}
Thus household support is both individually tapered and aggregately budgeted.

\begin{proposition}[Household bounds]
For every household $f$,
\begin{equation}
  0\le H_{f,t}\le h_t+s_t\left(K^*+\frac{\lambda}{1-\lambda}\right),
\end{equation}
and aggregate expenditure satisfies $0\le H_t\le\overline H_t$.
\end{proposition}

\begin{proof}
For $K\le K^*$, $\omega(K)=K\le K^*$. For $K>K^*$,
\[
\omega(K)=K^*+\sum_{m=1}^{K-K^*}\lambda^m
\le K^*+\frac{\lambda}{1-\lambda}.
\]
Since $0\le\xi_t\le1$, the individual bound follows. The definition of $\xi_t$ gives
$H_t=\min\{\widehat H_t,\overline H_t\}$ when $\widehat H_t>0$, proving the aggregate
bound.
\end{proof}

Because the base term $h_t$ could create incentives to split one economic household
into multiple administrative units, eligibility must rely on a legally defined care and
residence unit, with due-process protections for genuine family separation.

\section{Residual pools and exact allocation}

After mandatory claims and actual household support, the remaining capacity is
\begin{equation}
  Q_t=C_t-M_t-H_t\ge0.
\end{equation}
Choose policy shares $\theta_G,\theta_T,\theta_V\ge0$ satisfying
\begin{equation}
  \theta_G+\theta_T+\theta_V\le1.
\end{equation}
The authorized pools are
\begin{equation}
  \overline G_t^{+}=\theta_GQ_t,
  \qquad
  \overline T_t=\theta_TQ_t,
  \qquad
  \overline V_t=\theta_VQ_t.
\end{equation}

\subsection{Equal newborn bonus}

The cohort bonus is
\begin{equation}
  b_t^{+}=
  \begin{cases}
    \overline G_t^{+}/B_t, & B_t>0,\\
    0, & B_t=0.
  \end{cases}
\end{equation}
Every newborn $c\in\mathcal B_t$ receives
\begin{equation}
  g_{c,t}^{(0)}=\underline g_t+b_t^{+}.
\end{equation}
Consequently,
\begin{equation}
  \sum_{c\in\mathcal B_t}g_{c,t}^{(0)}
  =G_t^{\min}+\overline G_t^{+}
\end{equation}
when $B_t>0$. The realized cohort, rather than smoothed or forecast births, is the final
denominator. Forecasts may be used for budgeting, but never to calculate an entitlement
that can overspend its pool.

\subsection{Progressive account top-ups}

Let $y_{c,t}\ge0$ be an administratively defined measure of the resources available to
child $c$, preferably averaged across several years and adjusted for household size.
Let $y_t^*>0$ be a reference threshold. Define the bounded shortfall score
\begin{equation}
  z_{c,t}=\clamp{\frac{y_t^*-y_{c,t}}{y_t^*}}.
\end{equation}
Additional bounded indicators may represent disability, care burden, displacement, or
regional cost. For clarity, combine them into $d_{c,t}\in[0,1]$ and define
\begin{equation}
  a_{c,t}=1+\kappa_yz_{c,t}+\kappa_dd_{c,t},
  \qquad \kappa_y,\kappa_d\ge0.
\end{equation}
If $\mathcal C_t$ is nonempty, allocate
\begin{equation}
  \tau_{c,t}=\overline T_t
  \frac{a_{c,t}}{\sum_{j\in\mathcal C_t}a_{j,t}}.
\end{equation}
If $\mathcal C_t$ is empty, set every $\tau_{c,t}=0$ and $T_t=0$. Otherwise set
$T_t=\sum_c\tau_{c,t}$.

\begin{proposition}[Exact pool exhaustion]
If $\mathcal C_t\ne\varnothing$, then $\tau_{c,t}>0$ for every eligible child and
\begin{equation}
  \sum_{c\in\mathcal C_t}\tau_{c,t}=\overline T_t.
\end{equation}
No income report or demographic realization can cause top-ups to exceed their pool.
\end{proposition}

\begin{proof}
Since $a_{c,t}\ge1$, every denominator and every individual weight is positive. Summing
the normalized shares gives one.
\end{proof}

The optional adult distribution satisfies $0\le V_t\le\overline V_t$. Unallocated
capacity is not automatically spent.

\section{Annual budget closure and fund accounting}

Actual program distribution is
\begin{equation}
  D_t=U_t+G_t^{\min}+H_t+G_t^{+}+T_t+V_t,
\end{equation}
where $G_t^{+}=\overline G_t^{+}$ when $B_t>0$ and $0$ otherwise.

Current distributable resources, including a backstop if activated, are
\begin{equation}
  C_t^Y=(1-\phi_t)Y_t+Z_t.
\end{equation}
Actual source amounts are defined after actual spending:
\begin{align}
  D_t^Y&=\min\{D_t,C_t^Y\},\\
  D_t^W&=D_t-D_t^Y,\\
  E_t&=C_t^Y-D_t^Y.
\end{align}
Here $E_t\ge0$ is unused current revenue, which is transferred to the fund rather than
silently disappearing from the accounts. The fund evolves as
\begin{equation}\label{eq:fundlaw}
  W_{t+1}=W_t(1+r_t)+\phi_tY_t+E_t-D_t^W.
\end{equation}

\begin{proposition}[Budget closure]
Under the allocation rules above,
\begin{equation}
  0\le D_t\le C_t,
  \qquad
  0\le D_t^W\le\overline d_t.
\end{equation}
Moreover, the accounting identity
\begin{equation}
  W_{t+1}=W_t(1+r_t)+Y_t+Z_t-D_t
\end{equation}
holds.
\end{proposition}

\begin{proof}
Mandatory spending satisfies $M_t\le C_t$. Household spending satisfies
$H_t\le\beta_H(C_t-M_t)$. The three residual pools sum to at most $Q_t$, so total
spending cannot exceed $C_t$. Since $C_t=C_t^Y+\overline d_t$, the residual definition
$D_t^W=\pos{D_t-C_t^Y}$ implies $D_t^W\le\overline d_t$. Finally,
$E_t=C_t^Y-D_t^Y$ and $D_t=D_t^Y+D_t^W$. Substitution into
Equation~\eqref{eq:fundlaw} gives
\[
W_{t+1}=W_t(1+r_t)+\phi_tY_t+C_t^Y-D_t
=W_t(1+r_t)+Y_t+Z_t-D_t.
\]
\end{proof}

\begin{corollary}[Positive and family-independent base]
For every $c\in\mathcal B_t$,
\begin{equation}
  g_{c,t}^{(0)}\ge\underline g_t>0,
\end{equation}
and this lower bound is independent of $K_{f,t}$, parental income, and education-module
completion.
\end{corollary}

\section{Child-account dynamics}

An SSCF account is the legal property of the child, administered in trust until staged
access. It is not an asset of the parent, household, or general treasury. If accounts
are invested inside the social fund, the public balance sheet must record an equal
liability; otherwise the same assets would be counted both as public fund wealth and as
private child wealth.

Let $A_{c,s}$ be the account balance at the start of date $s$, $q_{c,s}>-1$ the net
credited return, $\tau_{c,s}\ge0$ the top-up, and $x_{c,s}\ge0$ the total amount paid out
at the end of the period. The law of motion is
\begin{equation}\label{eq:accountlaw}
  A_{c,s+1}=A_{c,s}(1+q_{c,s})+\tau_{c,s}-x_{c,s},
\end{equation}
subject to
\begin{equation}
  0\le x_{c,s}\le A_{c,s}(1+q_{c,s})+\tau_{c,s}.
\end{equation}
This equation records withdrawals explicitly. Without $x_{c,s}$, a staged-release rule
would repeatedly distribute claims against capital that remained unchanged on paper.

The public fund return $r_t$ and child-account return $q_{c,s}$ are distinct accounting
objects. If accounts are externally custodied, $q_{c,s}$ is simply their net portfolio
return. If internally custodied, the fund must segregate corresponding assets and
liabilities.

\section{Staged access without forfeiture}

Let release ages satisfy $a_1<a_2<a_3$, for example 18, 21, and 25. Choose target
shares $\eta_1,\eta_2,\eta_3\ge0$ with $\eta_1+\eta_2+\eta_3=1$. To avoid applying
fractions repeatedly to an already reduced balance, define conditional release rates
\begin{equation}
  k_j=
  \begin{cases}
  \displaystyle\frac{\eta_j}{1-\sum_{m<j}\eta_m}, & j=1,2,\\[8pt]
  1, & j=3.
  \end{cases}
\end{equation}

Let $P_{c,j}^{-}$ be the protected balance immediately before release $j$. The amount
that vests out of protection is
\begin{equation}
  R_{c,j}=k_jP_{c,j}^{-},
  \qquad
  P_{c,j}^{+}=P_{c,j}^{-}-R_{c,j}.
\end{equation}
Thus $k_1=eta_1$, $k_2=eta_2/(1-\eta_1)$, and the final release transfers all
remaining protected assets. Returns earned between milestones remain attached
proportionally to the protected assets that produced them.

Financial education affects the form of early access, not ownership. Let
\begin{equation}
  e_{c,j}=\sum_{m=1}^{J}\beta_mz_{c,m,j},
  \qquad z_{c,m,j}\in\{0,1\},
  \quad \beta_m\ge0,
  \quad \sum_m\beta_m=1.
\end{equation}
With guaranteed discretionary fraction $\delta\in(0,1]$, the immediately unrestricted
part of release $j<3$ is
\begin{equation}
  R_{c,j}^{\mathrm{free}}
  =R_{c,j}\bigl[\delta+(1-\delta)e_{c,j}\bigr].
\end{equation}
The remainder is placed in a restricted wallet usable for a broad statutory list that
may include education, housing, mobility, tools, healthcare, debt avoidance, and
enterprise. It is not forfeited. At $a_3$, all remaining restrictions expire regardless
of module completion.

The education program must be free, accessible, multilingual, disability-accommodating,
and subject to appeal. The modules should be short and practical. They cannot become a
test of cultural conformity or an administrative device for permanent exclusion.

\begin{proposition}[No duplicate release]
Ignoring payouts from already released wallets, the staged rule transfers the entire
protected balance exactly once by age $a_3$.
\end{proposition}

\begin{proof}
The first release leaves fraction $1-\eta_1$. The second conditional rate transfers
$\eta_2/(1-\eta_1)$ of that protected balance, corresponding to the second target share
after accounting for the first release. The third release rate equals one and therefore
transfers the entire remaining protected balance. Each transfer is subtracted from
$P_{c,j}^{-}$ before the next period.
\end{proof}

\section{Real per-capita sustainability}

Nominal fund non-decline is not an adequate definition of intergenerational
sustainability. A fund can grow in currency units while shrinking after inflation or
relative to the population it serves.

Define inflation and population growth by
\begin{equation}
  1+\pi_t=\frac{P_{t+1}}{P_t},
  \qquad
  1+n_t=\frac{N_{t+1}}{N_t},
\end{equation}
where $\pi_t>-1$ and $n_t>-1$. Real per-capita fund wealth is
\begin{equation}
  w_t=\frac{W_t}{P_tN_t}.
\end{equation}

\begin{proposition}[Exact real per-capita condition]
For $W_t>0$, real per-capita fund wealth is non-decreasing, $w_{t+1}\ge w_t$, if and only
if
\begin{equation}\label{eq:sustainability}
  1+r_t+\frac{\phi_tY_t+E_t-D_t^W}{W_t}
  \ge(1+\pi_t)(1+n_t).
\end{equation}
Equivalently,
\begin{equation}
  r_t+\frac{\phi_tY_t+E_t-D_t^W}{W_t}
  \ge \pi_t+n_t+\pi_tn_t.
\end{equation}
\end{proposition}

\begin{proof}
By the fund law,
\[
\frac{W_{t+1}}{W_t}
=1+r_t+\frac{\phi_tY_t+E_t-D_t^W}{W_t}.
\]
Also,
\[
\frac{P_{t+1}N_{t+1}}{P_tN_t}=(1+\pi_t)(1+n_t).
\]
The condition $W_{t+1}/(P_{t+1}N_{t+1})\ge W_t/(P_tN_t)$ is therefore equivalent to
Equation~\eqref{eq:sustainability}.
\end{proof}

This is an ex-post accounting condition, not a promise that risky returns will realize.
Ex-ante policy should use conservative expected real returns, rolling asset values,
stress tests, and a stabilization reserve. A draw cap $\gamma_t$ is a prudential tool;
it is not, by itself, proof of sustainability.

Where the relevant liability is more closely related to child cohorts than to total
population, the same proposition may be stated using a transparent demographic index
$L_t$ in place of $N_t$. The denominator should match the right being stabilized. Birth
counts determine cohort bonuses directly; total population growth should not be used as
an indirect substitute.

\part{Calibration, Governance, and Evaluation}

\section{Illustrative calibration}

The following values are arbitrary currency units (CU), not estimates for a particular
country. They demonstrate accounting behavior.

Assume
\begin{equation*}
  Y_t=100\text{ billion},\quad
  W_t=500\text{ billion},\quad
  \phi_t=0.40,\quad
  \gamma_t=0.03,\quad
  \widetilde W_t=W_t,
\end{equation*}
and $r_t=0.05$. Then
\begin{equation*}
  \overline d_t=15\text{ billion},
  \qquad
  C_t^0=60+15=75\text{ billion}.
\end{equation*}

Suppose
\begin{equation*}
  U_t=25\text{ billion},\qquad
  B_t=500{,}000,qquad
  \underline g_t=10{,}000.
\end{equation*}
The aggregate statutory base is $5$ billion and mandatory expenditure is
\begin{equation*}
  M_t=25+5=30\text{ billion}.
\end{equation*}
No backstop is required, so $Z_t=0$ and $C_t=75$ billion. Capacity after mandatory
claims is $45$ billion.

Let $\beta_H=0.50$. The household-layer budget is $22.5$ billion. If desired aggregate
household support equals $20$ billion, it is paid in full: $H_t=20$ billion. Remaining
capacity is
\begin{equation*}
  Q_t=75-30-20=25\text{ billion}.
\end{equation*}
Choose
\begin{equation*}
  \theta_G=0.60,\qquad
  \theta_T=0.20,\qquad
  \theta_V=0.
\end{equation*}
The newborn bonus pool is $15$ billion, so every newborn receives a bonus of
\begin{equation*}
  b_t^+=\frac{15{,}000{,}000{,}000}{500{,}000}=30{,}000.
\end{equation*}
The initial account is therefore $40{,}000$ per newborn: a $10{,}000$ protected base
plus a $30{,}000$ equal cohort bonus. The progressive top-up pool for all eligible child
accounts is $5$ billion. Total distribution is
\begin{equation*}
  D_t=25+5+20+15+5=70\text{ billion}.
\end{equation*}
Current distributable revenue is $60$ billion, so
\begin{equation*}
  D_t^Y=60\text{ billion},qquad
  D_t^W=10\text{ billion}.
\end{equation*}
The maximum draw is not taken merely because it is available. The closing fund is
\begin{equation*}
  W_{t+1}=500(1.05)+40-10=555\text{ billion}.
\end{equation*}

\subsection{Household-level illustration}

Let $h_t=5{,}000$, $s_t=4{,}000$, $K^*=2$, and $\lambda=0.60$. Before any aggregate
scaling,
\begin{align*}
  \omega(1)&=1,\\
  \omega(3)&=2.6,\\
  \omega(6)&=3.3056.
\end{align*}
Desired household supplements are therefore
\begin{align*}
  \widehat H(1)&=9{,}000,\\
  \widehat H(3)&=15{,}400,\\
  \widehat H(6)&=18{,}222.40.
\end{align*}
The taper bounds household support, while all newborns still receive the same cohort
capital.

\subsection{Account growth}

With no later top-ups or withdrawals and a constant net return of 4 percent, a $40{,}000$
initial account becomes
\begin{equation*}
  40{,}000(1.04)^{18}\approx81{,}032.66
\end{equation*}
at age 18. This is a nominal illustration; actual policy must report inflation-adjusted
values, fees, risk, and the distribution of realized returns.

\section{Demographic and revenue stress tests}

\subsection{Birth shock}

Suppose realized births rise from $500{,}000$ to $650{,}000$, while the child floor,
capacity, statutory base, and desired aggregate household support remain as above. The
base obligation rises to $6.5$ billion, so
\begin{equation*}
  M_t=31.5\text{ billion}.
\end{equation*}
Capacity after mandatory claims is $43.5$ billion. The household budget is $21.75$
billion, still enough to pay the $20$ billion desired amount. The new residual is $23.5$
billion. Hence
\begin{equation*}
  \overline G_t^+=0.60(23.5)=14.1\text{ billion},
\end{equation*}
and the per-newborn bonus becomes
\begin{equation*}
  b_t^+=\frac{14.1\text{ billion}}{650{,}000}
  \approx21{,}692.31.
\end{equation*}
Every child keeps the $10{,}000$ base, and total newborn expenditure remains exactly
within its authorized pool. The shock changes the bonus, not the legal minimum and not
the budget identity.

\subsection{Dedicated-revenue collapse}

Suppose instead that $C_t^0<M_t$. The backstop activates for exactly the shortfall
$Z_t=M_t-C_t^0$. Household support and all residual pools fall to zero if no capacity
remains. The system therefore protects essential childhood provision and the statutory
capital base before optional layers. This ordering makes the political tradeoff visible:
the state must either fund the promised right, amend it prospectively under valid legal
rules, or admit default. It cannot conceal insolvency by assigning a zero residual to a
nominally universal account.

\section{Comparative statics and design choices}

\subsection{The statutory base $\underline g_t$}

A higher base strengthens the independent child claim but raises mandatory expenditure
linearly:
\begin{equation*}
  \frac{\partial G_t^{\min}}{\partial\underline g_t}=B_t.
\end{equation*}
It also increases the probability that the fiscal backstop is needed. The base should
therefore be indexed and meaningful but calibrated to durable revenue, while variable
prosperity is shared through the cohort bonus.

\subsection{The household share $\beta_H$}

A larger $\beta_H$ protects more residual capacity for household support. When desired
support is below its cap, increasing $\beta_H$ has no effect. When the cap binds, it
raises $H_t$ and correspondingly reduces the residual available for bonuses and top-ups.

\subsection{The taper parameters $K^*$ and $\lambda$}

Increasing $K^*$ gives full marginal household weight to more children. Increasing
$\lambda$ slows the post-threshold taper. Neither parameter changes the child's
statutory or cohort capital grant.

\subsection{Saving and draw parameters}

A larger $\phi_t$ reduces current-flow capacity and raises fund saving. A larger
$\gamma_t$ increases maximum current capacity but does not require a larger draw.
Actual $D_t^W$ is determined by actual expenditure after current resources are used.
Both parameters should be evaluated against the real per-capita condition rather than a
nominal balance target.

\subsection{Progressive weights}

Increasing $\kappa_y$ or $\kappa_d$ reallocates a fixed top-up pool; it does not increase
aggregate spending. This is an advantage of normalization. Distributional choices are
not allowed to masquerade as financing assumptions.

\section{Governance and legal form}

The mathematics defines claims and constraints. It cannot supply trustworthy
institutions by itself.

\subsection{Fund governance}

The social wealth fund should have:

\begin{itemize}
  \item statutory independence and a fiduciary mandate;
  \item professional management separated from benefit adjudication;
  \item public reporting of fees, risk, benchmarks, and voting policy;
  \item independent audit and legislative oversight;
  \item conflict-of-interest, procurement, and anti-corruption rules;
  \item a published asset-liability statement distinguishing public assets from child-
  account liabilities; and
  \item an emergency rule specifying when, by whom, and for how long ordinary draw
  limits may be altered.
\end{itemize}

Ring-fencing must be legal rather than rhetorical. Political raids can occur through
direct withdrawal, mandated below-market investments, optimistic return assumptions,
or the transfer of unrelated liabilities onto the fund. Governance rules must address
all four.

\subsection{Protection of child ownership}

Accounts should be non-seizable by ordinary household creditors and unavailable as
collateral before release. Guardians should not be able to redirect the balance. The
law must specify beneficiary status in adoption, foster care, guardianship changes,
death before release, and family dissolution. Administrative appeals should be simple,
free, and independent.

\subsection{Eligibility and migration}

Eligibility choices are unavoidable. Citizenship-only rules can exclude long-term
resident children; immediate eligibility can create political objections concerning
short-term residence. The framework does not dictate a single rule, but it requires the
rule to be prospective, administrable, and consistent with equal treatment. Possible
designs include vesting by years of childhood residence, partial portability, and
reciprocal arrangements across jurisdictions.

\subsection{Privacy and data minimization}

The universal base requires little means testing. Progressive top-ups require more data
and therefore create privacy, stigma, and error costs. The state should use the smallest
resource measure necessary, publish the formula, limit retention, prohibit unrelated
use, and allow correction. Complexity should be justified by the additional equity it
produces.

\section{Macroeconomic and political-economy considerations}

A national program of meaningful scale can affect asset prices, labor supply, savings,
fertility, educational choices, migration, and credit markets. A grant usable for
housing may partly capitalize into land prices if housing supply is constrained. A
grant usable for tuition may be captured by providers if price regulation and supply do
not respond. Capital access can increase entrepreneurship, but failure risk and market
power shape who retains the gains.

The funding instruments also have different incidence and volatility. Resource rents
are geographically concentrated and cyclical. Carbon charges should decline if climate
policy succeeds. Land-value taxation is comparatively difficult to avoid but requires
competent valuation. Inheritance taxation raises questions of valuation, avoidance,
liquidity, family businesses, and cross-border mobility. The OECD concludes that well-
designed inheritance taxation can raise revenue and advance equality of opportunity,
while also emphasizing design tradeoffs and behavioral responses \cite{oecd2021}.

These differences argue against treating $Y_t$ as a single permanently stable stream.
Empirical implementation should model each component separately, estimate covariance
with fund returns and benefit needs, and stress-test severe joint shocks. A stabilization
reserve and conservative statutory base are preferable to optimistic promises.

Political economy matters as much as portfolio theory. Universal components can build
broad support, while progressive components direct more resources toward constrained
children. Universality without progressivity may leave inherited hierarchy largely
intact; progressivity without a common claim may become stigmatized and politically
fragile. The SSCF deliberately combines the two while accounting for them separately.

\section{Objections and replies}

\subsection{The framework penalizes saving or family provision}

The equal base is not withdrawn from children whose parents save. Progressive top-ups
may decline with household resources, but they are allocated from a distinct pool using
a gradual bounded score. This preserves a universal ownership relation while directing
additional help where capital constraints are greater.

\subsection{Parents, not the state, should provide for children}

Family duties and public duties are not mutually exclusive. Families already act inside
legal, monetary, educational, and infrastructural systems. The proposal does not
eliminate parental provision; it establishes a minimum social inheritance that cannot
be erased by parental poverty, neglect, death, or family size.

\subsection{Young adults will waste the grant}

Some recipients will make poor decisions, as do some private heirs. That fact does not
justify denying ownership only to people whose inheritance is publicly constituted.
Staged access, accessible education, anti-fraud enforcement, and broad restricted uses
can reduce predictable harms. Permanent paternalism would contradict the agency the
capital is meant to support.

\subsection{Public services are better than cash or capital}

Often they are better for particular purposes. That is why the development floor has
first priority. But services do not fully provide mobility, bargaining power, or the
ability to initiate a project. The correct comparison is not services \emph{or} capital;
it is a carefully ordered bundle.

\subsection{The common-inheritance category is politically manipulable}

It is. The response is a statutory taxonomy, non-overlap rules, transparent incidence
analysis, and democratic revision---not the fiction that existing property categories
are pre-political. The framework intentionally limits its dedicated base rather than
declaring all wealth socially claimable.

\subsection{A universal program cannot repair historical injustice}

Correct. It is not designed to settle all rectificatory claims. Universal common
inheritance, progressive capability policy, and specific rectification should be
debated on their own grounds and may operate simultaneously.

\section{Empirical evaluation agenda}

Implementation should begin with pilots or phased cohorts accompanied by pre-registered
evaluation. Relevant outcomes include:

\begin{itemize}
  \item childhood health, educational continuity, and housing stability;
  \item account accumulation, fees, returns, and loss dispersion;
  \item post-release education, mobility, home formation, enterprise, and debt;
  \item labor supply and job-search duration;
  \item fertility and household-formation responses;
  \item price capitalization in housing, tuition, and other eligible uses;
  \item take-up, administrative error, appeal outcomes, and distribution by protected
  characteristics;
  \item fund performance and real per-capita sustainability; and
  \item political durability and public understanding.
\end{itemize}

Evaluation should compare the full bundle with plausible alternatives: equivalent
spending on services, unrestricted transfers, debt reduction, or tax credits. The
question is not whether capital has benefits in isolation, but whether this architecture
outperforms realistic uses of the same fiscal resources.

\section{Limitations}

This paper supplies a normative architecture and an internally balanced accounting
model. It does not estimate behavioral elasticities, an optimal grant, or an optimal
portfolio. It does not resolve the appropriate jurisdiction, the incidence of every
funding instrument, or the interaction with existing welfare systems. It does not show
that the SSCF is superior in every country, particularly where basic state capacity or
essential services are absent.

The model also abstracts from uncertainty at the individual level. Returns differ,
households move, administrative data contain errors, and some children require costly
support not captured by a scalar adjustment. Aggregate budget closure is necessary but
not sufficient for justice.

Finally, capital is not power-neutral. Asset managers, approved-use providers, lenders,
and platform firms may attempt to capture the new balances. Consumer protection,
competition policy, housing supply, and public provision are complements to the SSCF,
not external details.

\section{Conclusion}

No generation creates its world from nothing. It inherits natural assets, accumulated
knowledge, public infrastructure, legal institutions, ecological liabilities, and a
distribution of property shaped by earlier decisions. Private effort and voluntary
exchange remain morally and economically important, but they do not exhaust the sources
of social wealth.

The Social Starting Capital Framework gives institutional form to that observation. It
does not promise equal outcomes. It guarantees a child-development floor and a positive
capital base; recognizes households without allowing family size to diminish the
child's independent claim; distributes bonuses and top-ups through exact pools; records
withdrawals and liabilities honestly; and judges fund sustainability in real per-capita
terms.

The proposal's moral test is simple: whether a child is treated as a future agent rather
than merely as a dependent member of a fortunate or unfortunate household. Its fiscal
test is equally simple: whether every promised unit is funded, every risk bearer is
named, and future cohorts are included in the accounting. A society that claims human
worth is not determined by birth should be willing to build institutions in which that
claim has material content.

\appendix

\section{Implementation sequence}

A cautious implementation can proceed in the following order:

\begin{enumerate}[label=\arabic*.]
  \item Define the universal development package and measure its incremental cost.
  \item Legislate the statutory base grant, index, eligibility rule, and fiscal backstop.
  \item Establish the fund's asset-liability boundary and independent governance.
  \item Dedicate non-overlapping revenue streams and publish incidence estimates.
  \item Build account custody, registration, guardianship, and appeal systems.
  \item Introduce household supplements under an aggregate cap.
  \item Add normalized progressive top-ups only when data quality and privacy safeguards
  are adequate.
  \item Pilot staged access and education rules, with all restrictions expiring at the
  final release age.
  \item Publish annual budget closure, stress tests, and the real per-capita metric.
  \item Evaluate outcomes against alternative uses of equal fiscal cost before scaling.
\end{enumerate}

\section{Compact algorithm}

For each fiscal year $t$:

\begin{enumerate}[label=\arabic*.]
  \item Observe $W_t$, $\widetilde W_t$, realized $Y_t$, registered $\mathcal C_t$, and
  finalized $\mathcal B_t$.
  \item Compute $\overline d_t$, $C_t^0$, $U_t$, $G_t^{\min}$, $M_t$, $Z_t$, and $C_t$.
  \item Compute desired household support, its aggregate cap, $\xi_t$, and actual $H_t$.
  \item Compute $Q_t$ and the three residual pools.
  \item Divide the newborn bonus by actual $B_t$; normalize top-up weights over actual
  eligible children.
  \item Calculate actual $D_t$, then $D_t^Y$, $D_t^W$, and $E_t$.
  \item Update the fund using Equation~\eqref{eq:fundlaw} and publish both sides of the
  accounting identity.
  \item Credit individual accounts, record later withdrawals using
  Equation~\eqref{eq:accountlaw}, and report segregated assets and liabilities.
  \item Evaluate Equation~\eqref{eq:sustainability} ex post and publish forward stress
  tests before setting the next year's parameters.
\end{enumerate}

\begin{thebibliography}{99}
\raggedright

\bibitem{ackerman1999}
Bruce Ackerman and Anne Alstott.
\newblock \emph{The Stakeholder Society}.
\newblock Yale University Press, New Haven, 1999.

\bibitem{atkinson2015}
Anthony B. Atkinson.
\newblock \emph{Inequality: What Can Be Done?}
\newblock Harvard University Press, Cambridge, MA, 2015.

\bibitem{bauluz2020}
Luis Bauluz, Yajna Govind, and Filip Novokmet.
\newblock Global land inequality.
\newblock World Inequality Lab Working Paper 2020/10, 2020.
\newblock \url{https://wid.world/document/global-land-inequality-world-inequality-lab-wp-2020-10/}.

\bibitem{george1879}
Henry George.
\newblock \emph{Progress and Poverty}.
\newblock D. Appleton and Company, New York, 1879.

\bibitem{hamilton2010}
Darrick Hamilton and William A. Darity Jr.
\newblock Can ``baby bonds'' eliminate the racial wealth gap in putative
post-racial America?
\newblock \emph{The Review of Black Political Economy}, 37(3--4):207--216, 2010.
\newblock \href{https://doi.org/10.1007/s12114-010-9063-1}{doi:10.1007/s12114-010-9063-1}.

\bibitem{locke1689}
John Locke.
\newblock \emph{Two Treatises of Government}.
\newblock Awnsham Churchill, London, 1689.

\bibitem{marx1867}
Karl Marx.
\newblock \emph{Capital: A Critique of Political Economy, Volume I}.
\newblock Otto Meissner, Hamburg, 1867.

\bibitem{norwayfund}
Norwegian Ministry of Finance.
\newblock The Government Pension Fund.
\newblock Government of Norway, accessed October 2026.
\newblock \url{https://www.regjeringen.no/en/topics/the-economy/the-government-pension-fund/id1441/}.

\bibitem{nozick1974}
Robert Nozick.
\newblock \emph{Anarchy, State, and Utopia}.
\newblock Basic Books, New York, 1974.

\bibitem{oecd2021}
Organisation for Economic Co-operation and Development.
\newblock \emph{Inheritance Taxation in OECD Countries}.
\newblock OECD Tax Policy Studies 28, OECD Publishing, Paris, 2021.
\newblock \href{https://doi.org/10.1787/e2879a7d-en}{doi:10.1787/e2879a7d-en}.

\bibitem{ostrom1990}
Elinor Ostrom.
\newblock \emph{Governing the Commons: The Evolution of Institutions for Collective
Action}.
\newblock Cambridge University Press, Cambridge, 1990.

\bibitem{paine1797}
Thomas Paine.
\newblock \emph{Agrarian Justice}.
\newblock Paris, 1797.

\bibitem{piketty2014}
Thomas Piketty.
\newblock \emph{Capital in the Twenty-First Century}.
\newblock Harvard University Press, Cambridge, MA, 2014.

\bibitem{polanyi1944}
Karl Polanyi.
\newblock \emph{The Great Transformation}.
\newblock Farrar \& Rinehart, New York, 1944.

\bibitem{rawls1971}
John Rawls.
\newblock \emph{A Theory of Justice}.
\newblock Harvard University Press, Cambridge, MA, 1971.

\bibitem{rousseau1755}
Jean-Jacques Rousseau.
\newblock \emph{Discourse on the Origin and Foundations of Inequality Among Men}.
\newblock Marc-Michel Rey, Amsterdam, 1755.

\bibitem{sen1999}
Amartya Sen.
\newblock \emph{Development as Freedom}.
\newblock Alfred A. Knopf, New York, 1999.

\bibitem{sherraden1991}
Michael Sherraden.
\newblock \emph{Assets and the Poor: A New American Welfare Policy}.
\newblock M. E. Sharpe, Armonk, NY, 1991.

\bibitem{vanderweide2024}
Roy van der Weide, Christoph Lakner, Daniel Gerszon Mahler, Ambar Narayan, and
Rakesh Gupta.
\newblock Intergenerational mobility around the world: A new database.
\newblock \emph{Journal of Development Economics}, 166:103167, 2024.
\newblock \href{https://doi.org/10.1016/j.jdeveco.2023.103167}{doi:10.1016/j.jdeveco.2023.103167}.

\bibitem{vanparijs1995}
Philippe Van Parijs.
\newblock \emph{Real Freedom for All: What (If Anything) Can Justify Capitalism?}
\newblock Oxford University Press, Oxford, 1995.

\bibitem{wir2022}
Lucas Chancel, Thomas Piketty, Emmanuel Saez, Gabriel Zucman, et al.
\newblock \emph{World Inequality Report 2022}.
\newblock World Inequality Lab, 2022.
\newblock \url{https://wir2022.wid.world/}.

\end{thebibliography}

\end{document}
